\lecture{6}{Thu 18 Sep 2025 12:33}{Covariant derivatives}

Furthering the topic of covariance, we introduce covariant derivatives. Tensors are automatically invariant under coordinate transformations, but derivatives are not usually tensors. We define a covariant derivative $\nabla _\mu $ to be a derivative that transforms in the right way. Since $\partial _\mu $ transforms properly on scalar functions, we should have $\nabla_\mu f=\partial _\mu f$. We should also have that if $S,T$ are tensors, then $\nabla _\mu (S+T)=\nabla _\mu S+\nabla _\mu T$ (linearity). Finally, it should satisfy the product rule with respect to the tensor product.

Let's see how this acts on a dual vector $\omega _\nu $. We should have a tensor $\nabla _\mu \omega _\nu =X_{\mu \nu } $. The only linear maps on dual vectors are derivatives $\partial _\mu $ and matrics $C_{\mu \nu } ^\alpha $. The most general thing we can have is thus
\[
	\nabla _\mu \omega _\nu =a \partial _\mu \omega _\nu -C_{\mu \nu } ^\alpha \omega _\alpha ,
\]
for some scalar $a$. The reason we use a minus sign is so the right side can later become $\Gamma _{\mu \nu } ^\alpha $. Condition 1 forces $a=1$ for scalars, and the product rule forces $a=1$ for vectors.

Now we consider it on an arbitrary vector $v^\alpha $. We must have
\[
	\nabla _\mu v^\alpha =\partial _\mu v^\alpha -B_{\mu \beta }^{\alpha } v^\beta 
.\]
We should find a relation between $B$ and $C$. Interestingly, we can already do this. Since $\omega _\alpha v^\alpha $ is a scalar,
\[
	\nabla _\mu (\omega _\alpha v^\alpha )=\partial _\mu (\omega _\alpha v^\alpha )=\omega _\alpha \partial _\mu v^\alpha +\left( \partial _\mu \omega _\alpha  \right) v^\alpha 
.\]
Alterntatively, via the product rule,
\[
	\nabla _\mu (\omega _\alpha v^\alpha )=(\nabla _\mu \omega _\alpha )v^\alpha +\omega _\alpha \nabla _\mu v^\alpha =\left( \partial _\mu \omega _\alpha -C_{\mu \alpha } ^\beta \omega _\beta \right) v^\alpha +\omega _\alpha \left( \partial _\mu v^\alpha -B_{\mu \beta } ^\alpha v^\beta  \right) 
.\]
These two expressions must be equal. We cancel the prtials to obtain
\[
	0=-C_{\mu \alpha } ^\beta \omega _\beta  v^\alpha +  B_{\mu \beta } ^\alpha \omega _\alpha v^\beta 
.\]
Intuitively, $B_{\mu \alpha } ^\beta =-C_{\mu \alpha } ^\beta $. To prove this, since $\alpha ,\beta $ are dummy indices, we can rename them as
\[
	C_{\mu \alpha } ^\beta \omega _\beta v^\alpha =B_{\mu \alpha } ^\beta  \omega _\beta v^\alpha 
.\]
Since this holds for \emph{every} dual vector and \emph{every} vector, this gives our relation. This can be done pointwise by choosing $\omega ,v$ to pick out individual entries. We now have
\[
	\nabla _\mu v^\alpha =\partial _\mu v^\alpha +C_{\mu \beta } ^\alpha  v^\beta 
.\]
The product rule can generalize to all tensors by taking the product over vectors and duals. Arbitrary tensors generalize by taking the derivative, then adding a term for each index, viewing upper indices as like vectors and lowers like duals.

We are basically correcting the derivative to make it transform like a tensor by adding all the other terms. Any tensor satisfying these properties is a covariant derivative. We specialize by choosing a nice covariant derivative for GR.

We want $C^\rho _{\mu \nu } =C_{\nu \mu } ^\rho $, meaning $C$ is symmetric in $\mu ,\nu $. This is nice for physics reasons, since we can just symmetrize it and antisymmetrize it. The antisymmeterization transforms like a tensor, and if that existed, then that would be an additional field. We have not found this extra field, so we demand it not have an antisymmetric component.

We also require that $\nabla _{\mu } g_{\alpha \beta } =0$. We call this metric compatibility. This will help develop the equations of motion from the equivalence principle. This fixes $\nabla _\mu $ as follows. We have
\[
	\nabla _\mu g_{\alpha \beta } =\partial _\mu  g_{\alpha \beta } -C_{\mu \alpha } ^\lambda  g_{\lambda \beta } -C_{\mu \beta} ^\lambda  g_{\alpha \lambda } 
.\]
We have
\[
	\partial _\mu g_{\alpha \beta } =C_{\mu \alpha } ^\lambda g_{\lambda \beta } +C_{\mu\beta } ^\lambda g_{\alpha \lambda } 
.\]
Since the partial of the metric is $\partial _\mu g_{\alpha \beta } =\Gamma _{\mu \alpha } ^\lambda g_{\lambda \beta } +\Gamma _{\mu \beta } ^\lambda g_{\alpha \lambda } $,  we just obtain
\[
	C_{\alpha \beta } ^\lambda  =\Gamma _{\alpha \beta } ^\lambda 
.\]

Here is an example. The metric for flat space in polar coordinates is
\[
	\mathrm{d} s^2 = -\mathrm{d} t^2+\mathrm{d} r^2+r^2\mathrm{d} \theta ^2
.\]
We thus have $g_{\mu \nu } =\operatorname{ dia g}(-1,1,r^2)$. Computing $\Gamma _{\theta \theta } ^r$ at a point:
\[
	\Gamma _{\theta \theta } ^r=\frac{1}{2} g^{r\rho } \left( 2 \partial _\theta g_{\rho \theta } -\partial _\rho g_{\theta \theta }  \right) =\frac{1}{2} g^{rr} \left( 2 \partial _\theta g_{r\theta } -\partial _rg_{\theta \theta }  \right) =\frac{1}{2} \frac{1}{r^2} \left( -\partial _r r^2 \right) =-r
.\]

